\chapter*{Conclusion}

Le modèle de données et la manière dont il est modélisé dans un système de gestion des collections dirige l'implémentation de nouvelles fonctionnalités applicatives. Dans le cadre de ce mémoire de recherche, nous avons mesuré la complexité de la mise en place de l'indexation automatisée par IA, alors même que le modèle de données adapté aux collections audiovisuelles (FRBR) pose des difficultés de navigation et de compréhension documentaires pour les utilisateurs. 
L'analyse de l'adaptation du modèle FRBR aux archives audiovisuelles a notamment permis de pointer sa verticalité, conçue pour les notices bibliographiques, qui écarte les spécificités et les différences de corpus vidéo. En effet, nombreuses sont les institutions qui conservent des archives documentaires, non diffusées, issues de films amateurs ou expérimentaux, auxquelles les niveaux du FRBR restent mal adaptés. Ainsi, l'utilisation du FRBR et la navigation entre ses différents niveaux dépend de la nature du corpus audiovisuel et est, de fait, hétérogène.  Dans ce contexte, le projet de la mise en place d'une indexation automatique des archives audiovisuelles multimodale pose question.

Dans le cadre du projet IA chez Skinsoft, nous avons tenté de proposer la mise en place d'un pipeline d'indexation automatique au sein de l'espace de travail S-Movies. Nous nous sommes toutefois confrontés à la complexité du reversement des métadonnées produites au sein du modèle FRBR, au delà des difficultés posées par un corpus audiovisuel historique et vieillissant. 

Dès lors, nous en avons conclu que la modélisation du FRBR selon le contexte institutionnel constitue un frein à la mise en place de l'indexation automatique, puisqu'elle repose sur des règles métier d'usage selon le corpus indexé. La refonte de la navigation entre les niveau du FRBR constitue alors un tournant essentiel et nécessaire à la mise en place d'une telle automatisation, permettant aux modèles d'apprentissage et à l'utilisateur de mieux comprendre les liens entre les différents niveaux d'indexation. 

Finalement, il s'agit de centrer notre approche sur l'expérience utilisateur en raison des difficultés d'entraîner un modèle d'intelligence artificielle adapté aux corpus audiovisuels et au modèle de données associé. Nous proposons alors la mise en place d'une aide à l'indexation, d'une indexation augmentée, qui suggèrerait à l'utilisateur des choix en les explicitant, sans pour autant automatiser l'ensemble du processus.  Par ses choix de validation, l'utilisateur entraînerait lui-même le modèle à la fois sur les corpus utilisés et à la fois sur l'utilisation du modèle FRBR. Une telle interface d'échange pourrait permettre de rendre l'indexation plus efficace et d'améliorer la compréhension des liens entre les niveaux d'indexation du FRBR. 

